\documentclass[11pt]{article}
\usepackage[margin=1in]{geometry}
\usepackage{enumitem}
% =============================================================================
% Preambles/header.tex — shared LaTeX/Beamer preamble for this project
%
% Use from a Beamer lecture by prepending:
%     \input{header}
% and compiling with TEXINPUTS=../Preambles:$TEXINPUTS (the /compile-latex
% skill does this automatically).
%
% PALETTE CONTRACT. Color names defined here MUST match the names in
% Quarto/theme-template.scss so Beamer and Quarto renderings use the same
% palette. The scripts/check-palette-sync.sh script enforces this.
%
% To customize: edit the HEX values below AND the matching $variable in
% Quarto/theme-template.scss, then run scripts/check-palette-sync.sh.
% =============================================================================

% -----------------------------------------------------------------------------
% Engine detection + encoding
%
% Our /compile-latex skill uses XeLaTeX, where inputenc is unnecessary and
% can produce encoding warnings. Only load inputenc under pdfTeX.
% -----------------------------------------------------------------------------
\usepackage{iftex}
\ifPDFTeX
  \usepackage[utf8]{inputenc}
\fi

\usepackage{amsmath, amssymb, amsthm}
\usepackage{booktabs}
\usepackage{graphicx}

% -----------------------------------------------------------------------------
% PALETTE — mirrors Quarto/theme-template.scss variable names.
% Load xcolor and define colors BEFORE anything that references them
% (notably hyperref, hypersetup, and the Beamer color assignments).
% -----------------------------------------------------------------------------
\usepackage{xcolor}

\definecolor{primary-blue}{HTML}{012169}      % headings, accents
\definecolor{primary-gold}{HTML}{B9975B}      % emphasis, borders
\definecolor{highlight-yellow}{HTML}{F2A900}  % markers, alerts
\definecolor{light-bg}{HTML}{E8EDF5}          % title / section backgrounds
\definecolor{jet}{HTML}{1A1A1A}               % body text

% Semantic colors (used in diagrams and annotations)
\definecolor{positive}{HTML}{15803D}          % good / observed / identified
\definecolor{negative}{HTML}{B91C1C}          % bad / problematic / bias
\definecolor{neutral}{HTML}{525252}           % reference / context

% Highlight accents (match .hi-* classes in the SCSS)
\definecolor{hi-slate}{HTML}{314F4F}
\definecolor{hi-green}{HTML}{2E8B57}
\definecolor{hi-red}{HTML}{C41E3A}

% Semantic aliases so skill templates and snippets can write
% \textcolor{accent}{...} without hard-coding "primary-blue".
\colorlet{accent}{primary-blue}
\colorlet{accent2}{primary-gold}

% -----------------------------------------------------------------------------
% Hyperref — loaded AFTER xcolor + palette so link colors resolve.
% -----------------------------------------------------------------------------
\usepackage{hyperref}
\hypersetup{colorlinks=true, linkcolor=primary-blue, urlcolor=primary-blue,
            citecolor=primary-blue, pdfborder={0 0 0}}

% -----------------------------------------------------------------------------
% Course code — set once per deck (after \input{header}, before \begin{document})
% via \coursecode{CS401}. Shown in the footer on every slide so a deck's course
% is identifiable without opening the title page. Empty by default.
% -----------------------------------------------------------------------------
\makeatletter
\newcommand{\coursecode}[1]{\gdef\@coursecodetext{#1}}
\gdef\@coursecodetext{}
\makeatother

% -----------------------------------------------------------------------------
% Beamer theme assignments (only apply under Beamer).
%
% We detect Beamer via \@ifclassloaded{beamer}, which is the documented,
% non-internal way. This must be wrapped in \makeatletter/\makeatother
% because \@ifclassloaded contains an @-catcode character.
% -----------------------------------------------------------------------------
\makeatletter
\@ifclassloaded{beamer}{%
  \usefonttheme{professionalfonts}
  \setbeamercolor{structure}{fg=primary-blue}
  \setbeamercolor{frametitle}{fg=primary-blue}
  \setbeamercolor{title}{fg=primary-blue}
  \setbeamercolor{subtitle}{fg=primary-gold}
  \setbeamercolor{author}{fg=primary-blue}
  \setbeamercolor{institute}{fg=neutral}
  \setbeamercolor{date}{fg=primary-gold}
  \setbeamercolor{itemize item}{fg=primary-blue}
  \setbeamercolor{itemize subitem}{fg=primary-gold}
  \setbeamercolor{alerted text}{fg=primary-gold}
  \setbeamercolor{block title}{fg=white, bg=primary-blue}
  \setbeamercolor{block body}{bg=light-bg}
  % Semantic block variants: block=definition/concept (blue), exampleblock=
  % worked example (green/positive), alertblock=key takeaway or caution (gold).
  % Keep to <=2 colored boxes per slide across ALL three variants (INV-7).
  \setbeamercolor{block title example}{fg=white, bg=positive}
  \setbeamercolor{block body example}{bg=light-bg}
  \setbeamercolor{block title alerted}{fg=white, bg=primary-gold}
  \setbeamercolor{block body alerted}{bg=light-bg}

  % Minimal footer: course code (left) + page X / N (right), no navigation clutter
  \setbeamertemplate{navigation symbols}{}
  \setbeamertemplate{footline}{%
    \color{neutral}\footnotesize\hspace{0.7em}\@coursecodetext\hfill
    \insertframenumber/\inserttotalframenumber\hspace{0.7em}\vspace{0.3em}}
}{}
\makeatother

% -----------------------------------------------------------------------------
% TikZ setup — libraries the snippet gallery uses
% -----------------------------------------------------------------------------
\usepackage{tikz}
\usetikzlibrary{arrows.meta, positioning, calc, decorations.pathreplacing,
                fit, shapes.geometric, backgrounds}

% Shared TikZ styles — snippets and hand-written diagrams can pick these up
% via `\begin{tikzpicture}[every node/.style={...}]` or by naming specific
% styles (e.g., `\node[dag-node] (X) at (0,0) {X};`).
\tikzset{
  % Node styles for causal DAGs and similar
  dag-node/.style = {
    draw=primary-blue, fill=light-bg, rounded corners=2pt,
    minimum width=1.2cm, minimum height=0.9cm, align=center, font=\small
  },
  decision-node/.style = {
    draw=primary-gold, fill=white, diamond, aspect=1.8,
    minimum width=1.8cm, minimum height=1.2cm, align=center, font=\small,
    inner sep=1pt
  },
  % Edge styles — observed vs counterfactual
  observed-edge/.style = {-{Latex[length=2.5mm]}, thick, draw=jet},
  counterfactual-edge/.style = {-{Latex[length=2.5mm]}, thick, draw=neutral, dashed},
  confound-edge/.style = {-{Latex[length=2.5mm]}, thick, draw=negative, dashed},
  % Data-point styles for plots
  observed-dot/.style = {fill=primary-blue, draw=primary-blue, circle, inner sep=1.2pt},
  counterfactual-dot/.style = {fill=white, draw=neutral, circle, inner sep=1.2pt}
}

% -----------------------------------------------------------------------------
% Convenience macros
% -----------------------------------------------------------------------------
\newcommand{\muted}[1]{\textcolor{neutral}{#1}}
\newcommand{\key}[1]{\textcolor{primary-gold}{\textbf{#1}}}
\newcommand{\good}[1]{\textcolor{positive}{#1}}
\newcommand{\bad}[1]{\textcolor{negative}{#1}}

% Standout frame macro: full-bleed dark background for section transitions.
% Beamer-only (uses \begin{frame}/\setbeamercolor/\usebeamerfont) -- do not
% call from an article-class file (e.g. Notes/<CODE>/*-notes.tex). \muted,
% \key, \good, \bad, and \coursecode above are plain xcolor/gdef and are
% safe to use from either class; verified when Notes/article-class support
% was added.
% Expand: \transitionslide{Your transition title}
\newcommand{\transitionslide}[1]{%
  \begingroup
  \setbeamercolor{background canvas}{bg=primary-blue}
  \begin{frame}[plain]
    \centering
    \vfill
    {\usebeamerfont{title}\color{white}\Large #1\par}
    \vfill
  \end{frame}
  \endgroup
}

% Section divider frame (Beamer-only), an alternative to \transitionslide with a
% darker two-line style: near-black (jet) background, a small white label line
% above a larger gold title, and a thin gold rule along the bottom edge.
% Expand: \sectiondivider{Section 2}{Pointers}
\newcommand{\sectiondivider}[2]{%
  \begingroup
  \setbeamercolor{background canvas}{bg=jet}
  \begin{frame}[plain]
    \vspace*{3.0cm}
    {\color{white}\ttfamily\large #1\par}
    \vspace{0.5cm}
    {\color{primary-gold}\LARGE #2\par}
    \begin{tikzpicture}[remember picture, overlay]
      \draw[primary-gold, line width=2.5pt]
        ($(current page.south west)+(0,0.1)$) -- ($(current page.south east)+(0,0.1)$);
    \end{tikzpicture}
  \end{frame}
  \endgroup
}

% -----------------------------------------------------------------------------
% Channel watermark — "Concepts That Click" (YouTube).
%
% Article-class files (CompetitiveExam Q/A sets, Notes, Assignments) get a
% light diagonal draft watermark on every page plus a centered footer naming
% the channel. Beamer decks get a subtle TikZ background watermark instead
% (the footline already carries the course code). Centralized here so every
% MCQ PDF is branded without per-file edits.
% -----------------------------------------------------------------------------
\makeatletter
\@ifclassloaded{article}{%
  \usepackage{draftwatermark}
  \SetWatermarkText{Concepts That Click}
  \SetWatermarkScale{1}
  \SetWatermarkFontSize{2cm}
  \SetWatermarkLightness{0.86}
  \SetWatermarkAngle{45}
  \usepackage{fancyhdr}
  \pagestyle{fancy}
  \fancyhf{}
  \fancyfoot[C]{\footnotesize\textcolor{neutral}{Concepts That Click~|~YouTube: \href{https://www.youtube.com/@concepts.thatclick}{@concepts.thatclick}}}
  \renewcommand{\headrulewidth}{0pt}
  \renewcommand{\footrulewidth}{0pt}
  \fancypagestyle{plain}{%
    \fancyhf{}%
    \fancyfoot[C]{\footnotesize\textcolor{neutral}{Concepts That Click~|~YouTube: \href{https://www.youtube.com/@concepts.thatclick}{@concepts.thatclick}}}%
    \renewcommand{\headrulewidth}{0pt}%
    \renewcommand{\footrulewidth}{0pt}%
  }
}{}
\@ifclassloaded{beamer}{%
  \setbeamertemplate{background}{%
    \begin{tikzpicture}[remember picture, overlay]
      \node[rotate=45, opacity=0.055, align=center]
        at (current page.center)
        {\Huge\textcolor{neutral}{Concepts That Click}};
    \end{tikzpicture}%
  }
}{}
\makeatother 
\coursecode{JKSSB}
\renewcommand{\thesection}{15.\arabic{section}}
\newtheorem{example}{Example}[section]
\newenvironment{definitionbox}[1]{%
  \par\noindent\textbf{Definition (#1).}\ \itshape
}{\par\vspace{0.3em}}
\title{Printers and Displays --- Detailed Lecture Notes}
\author{JKSSB Computer Awareness}
\date{\today}

\begin{document}
\maketitle

\section{Printers and hard-copy output}
A printer is an output device. It accepts data from the computer and records
it as \textbf{hard copy} --- output on a physical medium, normally paper.
Output that stays on a screen is called \textbf{soft copy}; a printer is the
device that turns soft copy into hard copy. Printers are described from three
angles: by the \textbf{mechanism} they use (dot matrix, inkjet, laser,
thermal, 3D), by whether they touch the paper (\textbf{impact} or
\textbf{non-impact}), and by the \textbf{unit printed} (a character, a line,
or a page).

\begin{definitionbox}{Printer}
An output device that records computer output as a permanent, human-readable
copy on paper or another medium.
\end{definitionbox}

\begin{center}
\begin{tabular}{p{0.26\textwidth}p{0.62\textwidth}}
\toprule
\textbf{Basis} & \textbf{Types} \\
\midrule
By mechanism & Dot matrix, inkjet, laser, thermal, 3D printer \\
By contact & Impact (strikes the paper) vs non-impact (no contact) \\
By unit printed & Character printer, line printer, page printer \\
\bottomrule
\end{tabular}
\end{center}

The same machine can carry several labels at once. A laser printer, for
example, is a non-impact printer, a page printer, and a laser-mechanism
printer. In an exam item, read which label is being asked for before choosing.

\section{Impact and non-impact printers}
An \textbf{impact printer} forms characters by physically striking an inked
ribbon against the paper. The contact makes these printers relatively noisy,
but because the type presses through the paper, they can produce
\textbf{multiple carbon copies} of multi-part forms. Common impact printers
are the dot matrix printer, the daisy wheel printer, and line printers such
as the drum printer and chain printer.

A \textbf{non-impact printer} forms characters without striking the paper.
These printers are quieter and generally give better print quality, but they
cannot make carbon copies. Common non-impact printers are the inkjet printer,
the laser printer, and the thermal printer.

\begin{definitionbox}{Impact printer}
A printer that strikes an inked ribbon against the paper to form a character.
\end{definitionbox}

\begin{definitionbox}{Non-impact printer}
A printer that forms characters without touching the paper with a striking
mechanism.
\end{definitionbox}

\begin{center}
\begin{tabular}{lll}
\toprule
\textbf{Point} & \textbf{Impact} & \textbf{Non-impact} \\
\midrule
Contact with paper & Yes (strikes) & No \\
Noise & Louder & Quieter \\
Carbon copies & Possible & Not possible \\
Examples & Dot matrix, daisy wheel, & Inkjet, laser, thermal \\
           & drum/chain line printers & \\
\bottomrule
\end{tabular}
\end{center}

\begin{example}
An item lists ``inkjet, laser, thermal, dot matrix'' and asks which is an
impact printer. Trace the test: which device strikes an inked ribbon against
the paper? Only the dot matrix printer does. The other three form the image
without contact, so they are non-impact.
\end{example}

\section{Impact printers in detail}
\subsection*{Dot matrix printer}
The dot matrix printer is an impact printer. Its print head holds rows of
tiny pins or hammers. The pins strike an inked ribbon to build each character
as a pattern of dots inside a matrix. It is noisy and its print quality is
low, but it is cheap to run and can print carbon copies and multi-part
invoices and receipts. For this reason it survives in billing and point-of-sale
work.

\subsection*{Daisy wheel printer}
A daisy wheel printer is an impact printer that uses a spinning wheel with
raised type characters, much like an electric typewriter. It produces good
text quality but cannot print graphics or images well.

\subsection*{Line printers}
A line printer prints an entire line of text in one operation instead of one
character at a time. A \textbf{drum printer} uses a rotating drum that carries
the full character set, with a hammer striking the paper wherever a character
is needed. A \textbf{chain (band) printer} moves a chain or band of type
characters past the hammers. Line printers are high-speed impact printers
associated with older mainframes and large batch printing.

\begin{definitionbox}{Line printer}
A printer that prints one full line of characters at a time; the drum printer
and chain printer are line printers.
\end{definitionbox}

\subsection*{Print units}
Printers are often grouped by the unit they print. A \textbf{character
printer} prints one character at a time (dot matrix, daisy wheel). A
\textbf{line printer} prints one line at a time (drum, chain). A \textbf{page
printer} prints one complete page at a time (laser and LED printers). The
unit grows from character to line to page.

\section{Inkjet printers}
An inkjet printer is a non-impact printer. Its print head sprays tiny
\textbf{droplets of liquid ink} through fine nozzles onto the page. Because
there is no striking, it is quieter than an impact printer, and it prints
good-quality full-colour images and photographs on many kinds of paper. Its
drawbacks are that it is slower per page than a laser printer, that wet ink
can \textbf{smudge}, and that its ink cartridges must be replaced when empty.
Inkjet printers use liquid ink, not toner.

\begin{definitionbox}{Inkjet printer}
A non-impact printer that forms characters and images by spraying fine
droplets of liquid ink onto the paper.
\end{definitionbox}

\section{Laser printers}
A laser printer is a non-impact printer that uses a \textbf{laser beam} and
\textbf{toner}, a fine dry powder, in a xerographic (electrostatic) process.
The laser draws the page image onto a charged drum; the toner sticks to the
charged areas; the toner is then transferred to the paper and fused in place
with heat and pressure. The result is fast, sharp, high-quality text, which
is why laser printers are common in offices. A laser printer is a
\textbf{page printer}: it prints a whole page at a time. Because it uses
toner rather than liquid ink, its output does not smudge when it gets wet.

\begin{definitionbox}{Laser printer}
A non-impact page printer that uses a laser beam to draw an image on a
charged drum and toner powder to form the printed page.
\end{definitionbox}

\begin{center}
\begin{tabular}{lll}
\toprule
\textbf{Point} & \textbf{Inkjet} & \textbf{Laser} \\
\midrule
Contact & Non-impact & Non-impact \\
Image formed by & Liquid ink droplets & Laser beam + toner \\
Unit printed & Character/line & Page printer \\
Speed & Slower & Faster \\
Strengths & Colour photos & Sharp text, offices \\
\bottomrule
\end{tabular}
\end{center}

\section{Thermal printers}
A thermal printer is a non-impact printer that uses \textbf{heat}. It heats
specially treated \textbf{thermal paper} (or a heat-sensitive ribbon) to form
characters. In its simplest form it needs no ink or toner cartridge, and it
operates silently. It is used for point-of-sale receipts, tickets, barcode
labels, and older fax machines. It is usually limited to a single colour.
Thermal, laser, and inkjet printers are all non-impact; dot matrix is impact.

\section{3D printers}
A 3D printer builds a solid three-dimensional object \textbf{layer by layer}
from a digital model. This process is called \textbf{additive
manufacturing}. It does not print ink onto paper; instead it adds material
such as plastic, resin, or metal until the object is complete. 3D printers
are used for prototypes, models, spare parts, and some medical implants. The
word ``printer'' is used, but the output is a physical object rather than a
document.

\section{Plotters}
A plotter is an output device that \textbf{draws} line graphics, usually with
one or more pens that move over the paper. It works with \textbf{vector}
drawings, so it is ideal for large engineering drawings, maps, and CAD plans,
and it handles sheet sizes that ordinary printers cannot. A plotter draws
lines; a printer lays down dots or characters. Both are output devices. A
plotter must not be confused with a scanner or a digitizer, which are input
devices.

\begin{example}
An item says, ``A plotter is an input device.'' This is wrong. A plotter
produces output by drawing lines on paper, so it is an \textbf{output}
device. The items that capture drawings as data --- the scanner and the
digitizer tablet --- are the input devices.
\end{example}

\section{Display terminology}
A monitor displays soft copy, and several exact terms describe it.

\begin{definitionbox}{Pixel}
The smallest addressable element of a display; short for ``picture element''.
Every image on the screen is built from pixels.
\end{definitionbox}

\begin{definitionbox}{Resolution}
The number of pixels a display can show, written as columns $\times$ rows,
for example $1920 \times 1080$. More pixels mean a sharper image.
\end{definitionbox}

The sharpness of a printed page is measured in \textbf{DPI}, which stands for
\textbf{Dots Per Inch}.

\begin{definitionbox}{Refresh rate}
How many times per second a display redraws the image. It is measured in
\textbf{hertz (Hz)}. A higher refresh rate (for example 120 Hz) gives
smoother motion and less flicker.
\end{definitionbox}

\begin{definitionbox}{Response time}
How long a pixel takes to change from one state to another. It is measured
in \textbf{milliseconds (ms)}. A lower response time means less motion blur
or ghosting.
\end{definitionbox}

\textbf{Brightness} is the amount of light a display emits, commonly measured
in \textbf{nits} (candela per square metre, cd/m$^2$). \textbf{Contrast} is
the difference between the brightest white and the darkest black the display
can produce. The \textbf{contrast ratio} compares the two, for example
$1000:1$; a higher ratio gives deeper blacks and more vivid images. In short:
brightness is how much light, and contrast is the range between light and
dark.

\begin{center}
\begin{tabular}{lll}
\toprule
\textbf{Display term} & \textbf{Meaning} & \textbf{Unit} \\
\midrule
Pixel & Smallest addressable element & --- \\
Resolution & Number of pixels shown & pixels (or DPI for print) \\
Refresh rate & Screen redraws per second & hertz (Hz) \\
Response time & Time for a pixel to change & milliseconds (ms) \\
Brightness & Light emitted by the screen & nits (cd/m$^2$) \\
Contrast ratio & Brightest white : darkest black & ratio \\
\bottomrule
\end{tabular}
\end{center}

\section{Exam retrieval and common confusions}
Six distinctions prevent most errors on this topic.
\begin{enumerate}
  \item \textbf{Impact vs non-impact}. Only a printer that \emph{strikes} an
    inked ribbon is impact. Dot matrix, daisy wheel, and drum/chain line
    printers are impact; inkjet, laser, and thermal are non-impact.
  \item \textbf{Dot matrix is impact}. It is easy to group it with inkjet and
    laser by ``modern-ness'', but it strikes the ribbon, so it is impact.
  \item \textbf{Inkjet vs laser}. Inkjet sprays \emph{liquid ink}; laser uses
    a \emph{laser beam and toner}. Do not swap the consumables.
  \item \textbf{Page vs line printer}. A laser printer is a page printer
    (one page at a time); a drum or chain printer is a line printer (one line
    at a time).
  \item \textbf{Printer vs plotter}. A plotter \emph{draws} vector line
    graphics for engineering drawings; it is an output device, not a printer
    of characters and not an input device.
  \item \textbf{Display units}. Refresh rate is in \emph{Hz} (higher is
    better); response time is in \emph{ms} (lower is better); brightness is
    in \emph{nits}; resolution is a pixel count.
\end{enumerate}

\begin{example}
An item asks for the difference between an inkjet and a laser printer. The
correct contrast is the image-forming method and consumable: the inkjet sprays
liquid ink droplets, while the laser uses a laser beam and toner powder.
Both are non-impact. A choice that calls the laser a character printer, or
says the inkjet uses toner, is wrong.
\end{example}

\subsection*{Exercises}
\begin{enumerate}[label=\arabic*.]
  \item Define an impact printer and a non-impact printer, and give two
    examples of each.
  \item Explain how a dot matrix printer forms a character, and state whether
    it is impact or non-impact.
  \item Distinguish an inkjet printer from a laser printer in terms of the
    image-forming method and the consumable used.
  \item What is a page printer? Name one example and contrast it with a line
    printer.
  \item Define pixel, resolution, refresh rate, response time, brightness,
    and contrast, giving the unit of measurement for each where one applies.
\end{enumerate}

\subsection*{Solutions}
\begingroup\small
\begin{enumerate}[label=\arabic*.]
  \item An impact printer forms characters by striking an inked ribbon
    against the paper (dot matrix, daisy wheel, drum/chain line printers). A
    non-impact printer forms characters without striking the paper (inkjet,
    laser, thermal).
  \item The dot matrix print head has rows of pins that strike an inked
    ribbon to build each character as a pattern of dots in a matrix. Because
    it strikes the ribbon, it is an impact printer.
  \item An inkjet printer sprays fine droplets of liquid ink through nozzles;
    a laser printer uses a laser beam to draw an image on a charged drum and
    toner powder to form the page. The inkjet consumable is liquid ink; the
    laser consumable is toner.
  \item A page printer prints a whole page at a time; a laser printer is an
    example. A line printer prints a whole line at a time (drum or chain
    printer), a smaller unit than a page.
  \item A pixel is the smallest addressable element of a display (no unit);
    resolution is the number of pixels shown, as columns $\times$ rows (pixel
    count, or DPI for print); refresh rate is the number of screen redraws per
    second (hertz, Hz); response time is the time for a pixel to change
    (milliseconds, ms); brightness is the light emitted by the screen (nits,
    cd/m$^2$); contrast is the difference between brightest white and darkest
    black, expressed as a contrast ratio.
\end{enumerate}
\endgroup
\end{document}
